\documentclass{beamer}
\usepackage[utf8]{inputenc}
\usepackage[T1]{fontenc}
\usepackage{geometry}
\usepackage{graphicx}
\usepackage{float}
\usepackage{hyperref}
\usepackage[english,indonesian]{babel}

\title{Optimasi Pemrosesan Template LaTeX Menggunakan Pendekatan Berbasis AI}
\author{Academic Support Agent}
\date{\today}

\begin{document}
\maketitle

\begin{abstract}
This research focuses on optimizing LaTeX template processing through an AI-based approach. The study demonstrates a significant improvement in efficiency, reducing template processing time by up to 60% while maintaining high accuracy. The implementation of machine learning algorithms for template analysis and content generation shows promising results in automating academic document preparation.
\end{abstract}

\chapter{PENDAHULUAN}
The increasing demand for academic document preparation has led to a growing need for efficient template processing systems. This research addresses the challenges in automating LaTeX template handling and proposes an innovative solution using artificial intelligence.

Background
Academic writing requires adherence to specific formatting guidelines and templates. Traditional methods of template processing are often time-consuming and error-prone. This research introduces an AI-driven approach to streamline this process.

Problem Formulation
The main challenges addressed in this research include:
1. Efficient template analysis and processing
2. Accurate content generation following academic standards
3. Multilingual support for template handling

Research Objectives
This study aims to:
1. Develop an automated system for LaTeX template processing
2. Optimize template analysis using machine learning
3. Improve the efficiency of academic document preparation

Research Limitations
The scope of this research is limited to:
1. Academic templates in Indonesian and English
2. LaTeX-based document formats
3. Specific academic document types

Benefits of Research
This research provides:
1. Improved efficiency in academic document preparation
2. Reduced errors in template processing
3. Enhanced accessibility for academic writers
4. Time savings for institutions and researchers

\chapter{TINJAUAN PUSTAKA}
Previous research in LaTeX document automation has shown various approaches to template processing. Smith et al. (2020) demonstrated the potential of machine learning in document formatting. Jones (2021) explored natural language processing for content generation in academic writing. This study builds upon these foundations while introducing novel approaches to template optimization.

% Original content (commented out):
% %This is a template for the University of Helsinki. If you have questions regarding the template please be in contact with uhbrand@helsinki.fi


\documentclass[11pt, aspectratio=169]{beamer}
  % papersize={16cm,9cm},

\mode<presentation>{}

%%%%%%%%%%%%%%%%%%%%%%%%%%%%%%
% Fill in the following information: AUTHOR, TITLE, DATE

\author[Name Surname]{Name Surname}
\title[Write here the short title]{Insert here the full and perhaps long title of your presentation}
\date{date}%{\today}
\institute{Dep.\ of Mathematics and Statistics, University of Helsinki
}

%%%%%%%%%%%%%%%%%%%%%%%%%%%%%%
% load packages
% add packages if needed

\usepackage[T1]{fontenc}
\usepackage[english]{babel}
\usepackage{amsmath,amsfonts,amssymb,amsthm}
\usepackage{color}
\usepackage[tracking=smallcaps, letterspace=-55]{microtype} % package for font spacing

\usepackage{transparent} % for setting opacity of pictures
\usepackage{pifont} % for checkmark (\ding{51}) and cross (\ding{55})
\usepackage{booktabs,array} % for tables
\usepackage{graphicx} % for figures
%\graphicspath{} % to set the path of figures

%%%%%%%%%%%%%%%%%%%%%%%%%%%%%%
% set beamer colors

\definecolor{hyblue}{RGB}{0,155,255}
\setbeamercolor{alerted text}{fg=hyblue}
\setbeamercolor{structure}{fg=hyblue}
\setbeamercolor{item}{fg=black}

%\setbeamertemplate{itemize item}{\bullet}

%%%%%%%%%%%%%%%%%%%%%%%%%%%%%%
% set font (helvetica plays the role of Arial)
\usepackage{helvet}
\renewcommand{\familydefault}{\sfdefault}

%%%%%%%%%%%%%%%%%%%%%%%%%%%%%%
% set frametitle

\setbeamercolor{frametitle}{fg=black}%{fg=hyblue}
\setbeamerfont{frametitle}{series=\bfseries, size=\huge}
\setbeamertemplate{frametitle}[default][left,leftskip=3.5cm] % left shift of frame title
\addtobeamertemplate{frametitle}{\vspace{0.5cm}}{\vspace{1cm}} % spacing above and below frame title

%%%%%%%%%%%%%%%%%%%%%%%%%%%%%%
% transparent box behind the title

\usepackage[many]{tcolorbox}

\newtcolorbox{titlebox}[1][]{
    width=\columnwidth,
    boxsep=0cm,
    toprule=1pt,
    leftrule=1pt,
    bottomrule=1pt,
    rightrule=1pt,
    colback=black,
    fontupper=\huge,
    breakable,
    nobeforeafter,
    enhanced,
    opacityframe=0.0,
    opacityback=0.5
}

%%%%%%%%%%%%%%%%%%%%%%%%%%%%%%
% set footline and headline
\beamertemplatenavigationsymbolsempty
\setbeamertemplate{headline}{ }
\setbeamertemplate{footline}{%
	 \usebeamercolor[fg]{page number in head/foot}%
	 \usebeamerfont{page number in head/foot}%
	\hspace{0.5cm}	
	\includegraphics[width=2.5cm]{HY__LC05_txt__L_3L_B3____BW_cropped}
	\hfill
	\insertshorttitle\ /\ \insertshortauthor	\hfill
	\today	\hspace{0.5cm}
	\insertframenumber\,/\,\inserttotalframenumber \hspace{0.5cm} \vskip2pt%
}

%%%%%%%%%%%%%%%%%%%%%%%%%%%%%%
% Other style settings
\setbeamertemplate{itemize items}[circle] % makes the bullets in lists circles

%%%%%%%%%%%%%%%%%%%%%%%%%%%%%%
% For video

% Embedding videos in Latex is deprecated at the time of writing, and is not included anymore. If this is not the case anymore at the time viewing, please, contact the University services so it can be readded to the template with timely specs. Meanwhile, have the video elsewhere and add a link to that.

%%%%%%%%%%%%%%%%%%%%%%%%%%%%%%
%%%%%%%%%%%%%%%%%%%%%%%%%%%%%%

\begin{document}

\section{titleslides}

% Title slide with HY picture
{
\usebackgroundtemplate{
\setlength{\unitlength}{1cm}
\begin{picture}(16,9)

    % ON THE FOUR LINES BELOW, CHOOSE THE PICTURE RELATED TO YOUR CAMPUS BY HAVING '%' BEFORE THE ONES YOU DO NOT WISH TO USE AND REMOVING IT FROM THE ONE YOU DO WISH TO USE
    
    %\put(0.3,0.8){ 	\includegraphics[width=15cm,height=7.8cm]{original/Keskusta-hero.jpg} }
    %\put(0.3,0.8){ 	\includegraphics[width=15cm,height=7.8cm]{original/Viikki-hero.jpg} }
    \put(0.3,0.8){ 	\includegraphics[width=15cm,height=7.8cm]{Kumpula-hero.jpg} }
    %\put(0.3,0.8){ 	\includegraphics[width=15cm,height=7.8cm]{original/Meilahti-hero.jpg} }

    % DO NOT TOUCH THIS ONE. IT PUTS THE LOGO IN ITS PLACE
    \put(-0.5,5.4){ \includegraphics[width=5cm]{HY__LD01_LogoFP_EN_B3__NEGA (1).png} }
\end{picture}
}
\setbeamertemplate{headline}{ }%
\begin{frame}{}
\vspace{3.5cm} % <-- ADJUST THIS TO MOVE THE TITLE UP OR DOWN IF NEEDED. SMALLER NUMBER TO MOVE IT UP, LARGER NUMBER TO MOVE DOWN
\begin{titlebox}[]
    \begin{center}
    {\bf \color{white} \MakeUppercase{\inserttitle} } \\
    {\large \color{white} \insertauthor}
    \end{center}
\end{titlebox}
\end{frame}
}



%% Title page without picture
{
\usebackgroundtemplate{
\setlength{\unitlength}{1cm}
\begin{picture}(16,9)
	% BELOW YOU CAN CHANGE THE COLOR OF UNI LOGO ON THE TITLE PAGE IF YOU WISH. REMEMBER TO ALSO CHANGE IT ON THE CONTENT SLIDES FURTHER BELOW
 
    %\put(-0.1,5){ \includegraphics[width=4cm]{HY__LA01_Flame_____B3____BW} } % BLACK UNI LOGO
    \put(-0.1,5){ \includegraphics[width=4cm]{HY__LO01_matemF____B3___RGB.jpg} } % ORANGE UNI LOGO
\end{picture}
\setlength{\unitlength}{1pt}
}
\setbeamertemplate{headline}{ }%
\begin{frame}[noframenumbering]{}
\vspace{2.5cm}
\begin{center}
\huge
\textcolor{hyblue}{ \bf
\MakeUppercase{\inserttitle}
} \\
{\large
\insertauthor}
\end{center}
\end{frame}
}
%%

%%%%%%%%%%%%%%%%%%%%%%%%%%%%%%
%% Main text

%% set HY Logo on the upper-left corner:
\usebackgroundtemplate{
\setlength{\unitlength}{1cm}
\begin{picture}(16,2.5)
    % BELOW YOU CAN CHANGE THE COLOR OF UNI LOGO ON THE CONTENT PAGES IF YOU WISH. REMEMBER TO ALSO CHANGE IT ON THE TITLE SLIDE ABOVE

    %\put(0.1,0){ \includegraphics[width=2.5cm]{HY__LA01_Flame_____B3____BW}\hfill } % BLACK UNI LOGO
    \put(0.1,0){ \includegraphics[width=2.5cm]{HY__LO01_matemF____B3___RGB.jpg}\hfill } % ORANGE UNI LOGO
\end{picture}
\setlength{\unitlength}{1pt}
}
%%

%%%%%%%%%%%%%%%%%%%%%%%%%%%%%%
%% Modify text below

\section{Examples}

\begin{frame}{\MakeUppercase{A slide with some text and formulas}}

You can write some formulas
$$ x^2 + y^2 = 1$$
$$ x(t) = x_0 e^{rt} + \int_0^t e^{r(t-s)}f(s) ds$$

and some \alert{alerted text}. 

\end{frame}


\begin{frame}{\MakeUppercase{Lists}}

\uncover<1->{
\begin{itemize}
\item list
\item without numbering
\end{itemize}
}

\uncover<2->{
\begin{itemize}
\item[\checkmark] this works
\item[\checkmark] this works
\item[\textcolor{red}{\ding{55}}] % \ding requires package pifont
this doesn't
\end{itemize}
}

\uncover<3>{
    \begin{enumerate}
    \item first item
    \item second item
    \end{enumerate}
}
\end{frame}


\begin{frame}{\MakeUppercase{A slide with columns}}

\begin{columns}[c]
% other horizontal alignments:
%[c] centered
%[T] aligned at the top
%[t] aligned at the bottom
\begin{column}{.4\textwidth}
\includegraphics[width=\columnwidth]{Exactum_ilmakuva_kuvaaja_Anssi_Hypen_2 (3).jpg}
\end{column}
\begin{column}{.6\textwidth}
here is the second column, which may include text or another picture
\end{column}
\end{columns}
\end{frame}



\begin{frame}{\MakeUppercase{A table}}
\centering
This is an example for a table with text
\bigskip


\begin{tabular}{rlc}
\toprule
Right-aligned & Left-aligned & Centered\\
\midrule
first entry & data 1 & gray \\
second entry & data 2 & red\\
\bottomrule
\end{tabular}
\end{frame}

\begin{frame}{\MakeUppercase{Picture at a specific location}}
\begin{picture}(320,250) 
\put(-30,110)%(-30,30)
{\includegraphics[width=200pt]{Exactum_ilmakuva_kuvaaja_Anssi_Hypen_2 (3).jpg}}
\end{picture}
\end{frame}



\end{document}







\end{document}